% --------------------------------------------------------------------------------
% Student guideline for the final year dissertation using LaTeX2e
% Alexandra Bonnici
% Department of Systems and Control
% 13th September 2017

\documentclass[12pt, oneside]{Thesis}

% --------------------------------------------------------------------------------
% Some useful packages
% You may add/remove any packages as you deem fit,
% Recommended packages which are recommended are marked by [R] in the comment
% --------------------------------------------------------------------------------

% font options & general typesetting
\usepackage{microtype}            % some refinements over the general appearance of text
\usepackage{eqparbox}             % ensures that blocks of text occupy the same space

% for figures
\usepackage[pdftex]{graphicx}     %  [R] to use graphics
\usepackage{subfigure}            %  [R] to make subfigures
\usepackage{epstopdf}             %  [R] adds eps support in the graphics package
\usepackage{wrapfig}              %  to allow text to wrap around figures


% for tables
\usepackage[table]{xcolor}
\usepackage{multirow,bigstrut}    % to have tables with merged rows
\usepackage{array, booktabs}      %  [R] for good looking tables
\usepackage{longtable}            % to create tables that span more than one page
\usepackage{adjustbox}            % so that long tables fit within the page
\usepackage{rotating}             % adds support to rotate pages

% to write nice mathematical equations
\usepackage{mathtools}            %  [R] some mathematical tools
\usepackage{amsmath}		      %  [R] to write math formulas easily	
\usepackage{amssymb}              %  [R] for maths symbols
\usepackage{amsthm}               % for the typesetting of theorems
\usepackage{esvect}               % to make vector arrows
\usepackage{bm}                   % to make bold maths symbols

\usepackage{xfrac}                % [R] to make in-text fractions

% more symbols
\usepackage{pifont}               % Postscript standard Symbol and Dingbats fonts
\usepackage{gensymb}              % generic symbols fonts 
\usepackage{stackrel}             % to stack things on top of each other

% to include urls
\usepackage{url}                  % to write urls

% to typeset pseudo code algorithms and source codes
\usepackage{algorithm}            % to type set algorithms as pseudo code
\usepackage{algpseudocode}        % more layout options for algorithm pseudo code
\usepackage{listings}             % to add non-formatted text e.g. source code

% to typeset theorems, lemmas and definitions
\newtheorem{theorem}{Theorem}[chapter]
\newtheorem{lemma}{Lemma}
\newtheorem{definition}{Definition}

% --------------------------------------------------------------------------------
% Information relevant to your dissertation. 
% Substitute with your own name, supervisor and title details.
%--------------------------------------------------------------------------------

\thesistitle{\LaTeX\ Template for the Final Year Project Dissertation} % title
\supervisor{Prof. Kenneth P. \textsc{Camilleri}}    % supervisor
\cosupervisor{}  									% co-supervisor 
\authors{Alexandra \textsc{Bonnici}}				% your name
\keywords{Student, Guidelines, FYP dissertation}    % keywords
\studentcode{12345M}  								% ID code
\department{Department of Systems and Control Engineering}  % department name
\Udegree{Bachelor of Engineering (Hons.)} 
\subject{Undergraduate Dissertation}          
\UNIVERSITY{UNIVERSITY OF MALTA}       
\faculty{Faculty of Engineering}                                      

\newif\ifHaveACosupervisor
% toggle the comment on the true/false lines as appropriate
%\HaveACosupervisortrue	
\HaveACosupervisorfalse

%--------------------------------------------------------------------------------
% Creates the PDF meta-data (for the electronic version of your dissertation). 
% There is no need to modify this information
%--------------------------------------------------------------------------------
\hypersetup{urlcolor=blue,%
		    linkcolor=blue,% 
            citecolor=blue,%
            colorlinks=true,% 
            hypertexnames=true}
            
\hypersetup{pdftitle={\ttitle}}
\hypersetup{pdfsubject=\subjectname}
\hypersetup{pdfauthor=\authornames}
\hypersetup{pdfkeywords=\keywordnames}

%--------------------------------------------------------------------------------
% Creates the title page (typesetting is taken care of and details are as above)
%--------------------------------------------------------------------------------
\title{\ttitle}

%--------------------------------------------------------------------------------
% Creates the glossaries pages and index pages
%--------------------------------------------------------------------------------
\makeglossaries
\makeindex

%--------------------------------------------------------------------------------
% The actual start of the dissertation
%--------------------------------------------------------------------------------
\begin{document}

\frontmatter            % sets page numbers of the first few pages to roman 
\setstretch{1.3}        % ensures a one-and-a-half line spacing
\fancyhead{}            % clears the headers
\rhead{\thepage}        % sets the top right header to the page number
\lhead{}                % ensures that the top left header remains empty
\pagestyle{fancy}     	% uses the "fancy" page style (adds the line at top)


%--------------------------------------------------------------------------------
% Prints out the title page
%--------------------------------------------------------------------------------
\input{./Chapters/TitlePage}

% -------------------------------------------------------------------------------
% Prints out the copyright page (you do not need to modify this)
%--------------------------------------------------------------------------------
\clearpage
\input{./Chapters/Copyright}

\clearpage
\input{./Chapters/DeclarationPage}

%--------------------------------------------------------------------------------
% Prints out the abstract 
% You will need to change the text in the file to reflect your own abstract
%--------------------------------------------------------------------------------
\clearpage
\input{./Chapters/Abstract}


%--------------------------------------------------------------------------------
% Prints out the acknowledgements page 
% You will need to change the text to reflect your own acknowledgements
%--------------------------------------------------------------------------------
\clearpage
\input{./Chapters/Acknowledgements}


%--------------------------------------------------------------------------------
% Prints out the list of contents, figures and tables 
% You don't have to modify any of this
%--------------------------------------------------------------------------------
\clearpage
\pagestyle{fancy}

\lhead{\emph{Contents}}
\tableofcontents

\lhead{\emph{List of Figures}}
\listoffigures

\lhead{\emph{List of Tables}}
\listoftables

%--------------------------------------------------------------------------------
% Prints out the list of acronyms 
% No changes will be necessary if used properly within text
%--------------------------------------------------------------------------------
\clearpage
\setglossarystyle{altlist}
\printglossary[type=\acronymtype, title=List of Acronyms, toctitle=List of Acronyms]

%--------------------------------------------------------------------------------
% Prints out the list of symbols 
% No changes will be necessary if used properly within text, comment if not relevant
%--------------------------------------------------------------------------------

\clearpage
\setglossarystyle{list}
\printglossary[title=List of Symbols,toctitle=List of Symbols]


%--------------------------------------------------------------------------------
% Prints out the chapter contents
%--------------------------------------------------------------------------------

\mainmatter
\pagestyle{fancy}

\clearpage
\chapter{Introduction}
\label{Chapt:Introduction}
\lhead{Chapter \thechapter. \emph{Introduction}}
\input{./Chapters/Introduction}

%--------------------------------------------------------------------------------
% Prints out any appendices (comment/delete if you do not need any)
%--------------------------------------------------------------------------------
\clearpage
\addtocontents{toc}{\vspace{2em}}
\appendix
\baselineskip=16pt

\input{./Chapters/AppendixA}

\addtocontents{toc}{\vspace{2em}} % Add a gap in the Contents, for aesthetics

\backmatter

%--------------------------------------------------------------------------------
% Prints out the bibliography
%--------------------------------------------------------------------------------

\label{Bibliography}

\lhead{\emph{Bibliography}}
\renewcommand{\bibname}{References}
\bibliographystyle{MyBibliographyStyle}
\bibliography{./Chapters/FYPBibliography}


\end{document}  